\section{Forward-Looking Discussion}
\label{sec:forward}

Past behavior around the IT boom is a poor template for AI, and our results say exactly where. Table~\ref{tab:diffs} lists the differences we think matter and what each predicts in the data. The first two are already visible: the sign of the junior-employment effect flipped between eras. The third and fourth explain why measured productivity may lag. The fifth is the one that would most change the investment case.

\begin{table}[t]
\caption{How generative AI differs from the 1995--2004 IT boom, and where each difference should show up.}
\label{tab:diffs}
\begin{singlespace}\small\centering
\begin{tabular}{p{0.04\linewidth}p{0.38\linewidth}p{0.47\linewidth}}
\toprule
 & Difference & Prediction and status \\
\midrule
D1 & IT automated routine middle-skill tasks; AI targets codified cognitive tasks in white-collar work. & Knowledge work flips from complement to substitute. \emph{Seen:} C2 sign flip. \\
D2 & IT rewarded computer-literate graduates; AI substitutes for the entry-level tasks that juniors learn on. & Junior--senior gap opens only in the AI era. \emph{Seen:} C2 and the kink test. \\
D3 & IT required heavy user-side capital and reorganization; AI is rented per use. & Faster adoption, less capital deepening at users. \emph{Not yet tested.} \\
D4 & Purchased AI services are intermediate inputs; user value added rises only if output rises by more. & Measured user productivity understates gains. \emph{Consistent with} C1 null and C3. \\
D5 & IT hardware was competitive; AI supply (chips, frontier models) is concentrated. & Producers may capture the gain through pricing, shrinking user margins over time. \emph{Watch:} intermediate-input costs in exposed sectors. \\
D6 & The Fed let 1990s productivity growth run; 2022--26 is a hiking cycle, with a further hike in September 2026. & Hiring freezes in rate-sensitive sectors mimic substitution; discount-rate effects dominate returns. \emph{Seen:} C4 rate loading. \\
\bottomrule
\end{tabular}
\end{singlespace}
\end{table}

\paragraph{Current conditions.} The FOMC raised the federal funds target to 3.75--4\% in September 2026 with inflation above target. Two consequences follow for this trade. First, rate-sensitive sectors are also the AI-exposed ones, so any long-short on exposure carries duration risk; our strategy's returns after ChatGPT load positively on the 2-year yield ($t=2.6$, Table~\ref{tab:strategy}). Second, tight policy is itself a reason for firms to freeze junior hiring, which makes the labor signal harder to read in real time. Our nowcast suggests the junior gap kept widening through 2026, by about 1.7 log points over the year to 2026Q3.

\paragraph{What would change our view.} Three developments would make the trade more attractive: (i) a measurable productivity acceleration in exposed user industries, which would confirm that the margin gains are durable; (ii) evidence that AI vendors are \emph{not} capturing the surplus (falling per-unit AI prices that do not raise user input costs); and (iii) an easing cycle, which would reduce the rate component of sector returns. A reversal of the junior decline as the 2026 graduating class is absorbed would weaken the substitution reading altogether.

\paragraph{The bubble question.} The IT boom delivered real productivity gains and a crash in the stocks of the firms that sold the technology. Our evidence suggests a version of the same split: the economic gain from AI so far appears in the margins of the industries that \emph{use} it, while equity attention sits with the firms that \emph{sell} it. An investor worried about an AI valuation bubble in producers can hold the users whose labor data show substitution as a partial hedge; the hedge has economic logic even though, as the next section shows, it has not earned excess returns on its own.
